%%%%%%%%%%%%%%%%%
% This is an sample CV template created using altacv.cls
% (v1.7.4, 30 July 2025) written by LianTze Lim (liantze@gmail.com). Compiles with pdfLaTeX, XeLaTeX and LuaLaTeX.
%
%% It may be distributed and/or modified under the
%% conditions of the LaTeX Project Public License, either version 1.3
%% of this license or (at your option) any later version.
%% The latest version of this license is in
%%    http://www.latex-project.org/lppl.txt
%% and version 1.3 or later is part of all distributions of LaTeX
%% version 2003/12/01 or later.
%%%%%%%%%%%%%%%%

%% Use the "normalphoto" option if you want a normal photo instead of cropped to a circle
% \documentclass[10pt,a4paper,withhyper,normalphoto]{altacv}

\documentclass[10pt,a4paper,withhyper]{altacv}
%% AltaCV uses the fontawesome5 and simpleicons packages.
%% See http://texdoc.net/pkg/fontawesome5 and http://texdoc.net/pkg/simpleicons for full list of symbols.

% Change the page layout if you need to
\geometry{left=1.25cm,right=1.25cm,top=1.5cm,bottom=1.5cm,columnsep=1.2cm}

% The paracol package lets you typeset columns of text in parallel
\usepackage{paracol}

% Change the font if you want to, depending on whether
% you're using pdflatex or xelatex/lualatex
% WHEN COMPILING WITH XELATEX PLEASE USE
% xelatex -shell-escape -output-driver="xdvipdfmx -z 0" sample.tex
\iftutex
  % If using xelatex or lualatex:
  \setmainfont{Roboto Slab}
  \setsansfont{Lato}
  \renewcommand{\familydefault}{\sfdefault}
\else
  % If using pdflatex:
  \usepackage[rm]{roboto}
  \usepackage[defaultsans]{lato}
  % \usepackage{sourcesanspro}
  \renewcommand{\familydefault}{\sfdefault}
\fi

% Change the colours if you want to
\definecolor{SlateGrey}{HTML}{2E2E2E}
\definecolor{LightGrey}{HTML}{666666}
\definecolor{DarkPastelRed}{HTML}{450808}
\definecolor{PastelRed}{HTML}{8F0D0D}
\definecolor{GoldenEarth}{HTML}{E7D192}
\definecolor{StrawberryCream}{HTML}{E86D75}
\definecolor{YellowOcher}{HTML}{4E1500}
\colorlet{name}{black}
\colorlet{tagline}{PastelRed}
\colorlet{heading}{DarkPastelRed}
\colorlet{headingrule}{YellowOcher}
\colorlet{subheading}{PastelRed}
\colorlet{accent}{PastelRed}
\colorlet{emphasis}{SlateGrey}
\colorlet{body}{LightGrey}

% Change some fonts, if necessary
\renewcommand{\namefont}{\Huge\rmfamily\bfseries}
\renewcommand{\personalinfofont}{\footnotesize}
\renewcommand{\cvsectionfont}{\LARGE\rmfamily\bfseries}
\renewcommand{\cvsubsectionfont}{\large\bfseries}


% Change the bullets for itemize and rating marker
% for \cvskill if you want to
\renewcommand{\cvItemMarker}{{\small\textbullet}}
\renewcommand{\cvRatingMarker}{\faCircle}
% ...and the markers for the date/location for \cvevent
% \renewcommand{\cvDateMarker}{\faCalendar*[regular]}
% \renewcommand{\cvLocationMarker}{\faMapMarker*}


% If your CV/résumé is in a language other than English,
% then you probably want to change these so that when you
% copy-paste from the PDF or run pdftotext, the location
% and date marker icons for \cvevent will paste as correct
% translations. For example Spanish:
% \renewcommand{\locationname}{Ubicación}
% \renewcommand{\datename}{Fecha}


%% Use (and optionally edit if necessary) this .tex if you
%% want to use an author-year reference style like APA(6)
%% for your publication list
% \input{pubs-authoryear}

%% Use (and optionally edit if necessary) this .tex if you
%% want an originally numerical reference style like IEEE
%% for your publication list
\input{pubs-num}

%% sample.bib contains your publications
\addbibresource{sample.bib}

\begin{document}
\name{Dominic Darrigo}
\tagline{Research Physicist}
%% You can add multiple photos on the left or right
\photoR {4.0 cm}{atom-png-29.png}
% \photoL{2.5cm}{Yacht_High,Suitcase_High}

\personalinfo{%
  % Not all of these are required!
  \email{ddarrigo.workspace@gmail.com}
  \phone{+1(716)-817-3899}

  %\homepage{www.homepage.com}
  % \twitter{@twitterhandle}
  % \xtwitter{@x-handle}
  \linkedin{www.linkedin.com/in/dominic-darrigo}
  \orcid{0009-0000-1890-2991}
  %% You can add your own arbitrary detail with
  %% \printinfo{symbol}{detail}[optional hyperlink prefix]
  % \printinfo{\faPaw}{Hey ho!}[https://example.com/]

  %% Or you can declare your own field with
  %% \NewInfoFiled{fieldname}{symbol}[optional hyperlink prefix] and use it:
  \NewInfoField{gitlab}{\faGitlab}[https://gitlab.com/]
  %\gitlab{your_id}
  %%
  %% For services and platforms like Mastodon where there isn't a
  %% straightforward relation between the user ID/nickname and the hyperlink,
  %% you can use \printinfo directly e.g.
  % \printinfo{\faMastodon}{@username@instace}[https://instance.url/@username]
  %% But if you absolutely want to create new dedicated info fields for
  %% such platforms, then use \NewInfoField* with a star:
  %\NewInfoField*{mastodon}{\faMastodon}
  %% then you can use \mastodon, with TWO arguments where the 2nd argument is
  %% the full hyperlink.
  %\mastodon{@username@instance}{https://instance.url/@username}
}

\makecvheader
%% Depending on your tastes, you may want to make fonts of itemize environments slightly smaller
% \AtBeginEnvironment{itemize}{\small}

%% Set the left/right column width ratio to 6:4.
\columnratio{0.6}

% Start a 2-column paracol. Both the left and right columns will automatically
% break across pages if things get too long.
\begin{paracol}{2}
\cvsection{Experience}
\cvevent{President}{of Society of Physics Students}{May 2025--Ongoing}{University at Buffalo}
\begin{itemize}
    \item Organized club meetings
    \item Lead advertiser of club
    \item Orchestrated guest speakers for the club 
    \item Lead tutor for SPS outreach tutoring
    \item Previously Vice President (September 2024--May 2025)
\end{itemize}

\divider

\cvevent{Researcher}{under Dr. Dusan Sarenac PhD}{September 2024 -- Ongoing}{University at Buffalo}
\begin{itemize}
\item Lead researcher on Spacial Neutron Modulator
\item Assistant researcher on various projects
\end{itemize}

\divider

\cvevent{Intern}{at National Center for Neutron Research}{Febuary 2025}{NIST Gaithersburg, MD}
\begin{itemize}
    \item Wrote technical documents on setup and operation of high-end camera and associated software to image X-rays.
\end{itemize}

\divider

\cvevent{Researcher}{under Dr. Sambandamurthy Ganapathy PhD}{January 2024 -- December 2024}{University at Buffalo}
\begin{itemize}
\item Assistant undergraduate researcher in  NbOx phase change relaxation Oscillators
\end{itemize}

\cvsection{Projects}

\cvevent{Spatial Neutron Modulator}{Department of Physics at UB}{2 Years}{}
\begin{itemize}
\item Design a device that can modulate the phase of neutrons to desired specification
\item Constructed large-scale prototype of device from electronic components and coded the microprocessor used to modulated the phase attenuation
\end{itemize}

\medskip

% \cvsection{A Day of My Life}

% % Adapted from @Jake's answer from http://tex.stackexchange.com/a/82729/226
% % \wheelchart{outer radius}{inner radius}{
% % comma-separated list of value/text width/color/detail}
% \wheelchart{1.5cm}{0.5cm}{%
%   6/8em/accent!30/{Sleep,\\beautiful sleep},
%   3/8em/accent!40/Hopeful novelist by night,
%   8/8em/accent!60/Daytime job,
%   2/10em/accent/Sports and relaxation,
%   5/6em/accent!20/Spending time with family
% }

% use ONLY \newpage if you want to force a page break for
% ONLY the current column
\newpage

\cvsection{Publications}

%% Specify your last name(s) and first name(s) as given in the .bib to automatically bold your own name in the publications list.
%% One caveat: You need to write \bibnamedelima where there's a space in your name for this to work properly; or write \bibnamedelimi if you use initials in the .bib
%% You can specify multiple names, especially if you have changed your name or if you need to highlight multiple authors.
\mynames{Darrigo\bibnamedelima Dominic \bibnamedelima Santo}
%% MAKE SURE THERE IS NO SPACE AFTER THE FINAL NAME IN YOUR \mynames LIST

\nocite{*}

\printbibliography[heading=pubtype,title={\printinfo{\faBook}{Patents}},type=patent]


%% Switch to the right column. This will now automatically move to the second
%% page if the content is too long.
\switchcolumn

\cvsection{My Life Philosophy}

\begin{quote}
``Never leave something only half done. Because things will never be only half broken when you need to fix them.''
\end{quote}

\cvsection{Most Proud of}

\cvachievement{\faUniversity}{Inducted Member of $\Sigma \Pi \Sigma$}{the honors society of the national Society of Physics Students, inductees are only nominated by professors in the SPS}

\divider
\cvachievement{\faUniversity}{Member of University of Buffalo Honors College}{Inducted member since enrollment.}
\divider


\cvachievement{\faLightbulb}{Author of Three Patents}{the patents detail the design and method of operations for vehicle emissions testing in a small form factor; suited for use in third-world countries}


\cvsection{Strengths}

% Don't overuse these \cvtag boxes — they're just eye-candies and not essential. If something doesn't fit on a single line, it probably works better as part of an itemized list (probably inlined itemized list), or just as a comma-separated list of strengths.

\cvtag{Hard-working}
\cvtag{Task-Oriented}\\
\cvtag{Driven by curiosity}
\cvtag{Motivator \& Leader}\\


\cvsection{Skills}

\cvtag{Complex Problem analysis}
\cvtag{Advanced Mathematics}
\cvtag{Trained in Nuclear Safety}\\
\cvtag{Author of Technical and Academic papers}\\
\cvtag{Knowledgeable in CAD Modeling}


\cvsection{Languages}

\cvskill{English}{5}
\divider

 %% Supports X.5 values.
\cvskill{LaTeX}{5}
\divider
\cvskill{Python}{3}
\divider
\cvskill{Mathematica}{4}
\divider
\cvskill{Arduio (C++)}{4.5}
%% Yeah I didn't spend too much time making all the
%% spacing consistent... sorry. Use \smallskip, \medskip,
%% \bigskip, \vspace etc to make adjustments.
\bigskip

\cvsection{Education}

\cvevent{B.Sc.\ in Physics}{University at Buffalo}{August 2022 -- May 2026}{}

\divider

\cvsection{Referees}

% \cvref{name}{email}{mailing address}
\cvref{Prof.\ Dusan Sarenac, PhD}{University at Buffalo}{dusansar@buffalo.edu}
{University at Buffalo\\
261 Fronczak Hall}

\divider

\cvref{Dr. Michael G. Huber}{NIST}{michael.huber@nist.gov}
{NIST Center for Neutron Research\\Gaithersburg, MD}


\end{paracol}


\end{document}
